% GENERATED FILE. Edit canonical lessons, then run npm run generate:books.
% canonical-source-hash: fnv1a64:80a9bceaccf095ac
% canonical-lessons: ML-C31-shubha-madhyaahnam, ML-S125-letter-uu, ML-C31-afternoon-convergence

\chapter{Good Afternoon}
\label{ch:ml-good-afternoon}
\hlchaptermodality{31}

\begin{chapteropening}
\textbf{By the end of this chapter:} \emph{I can greet someone in the afternoon in Malayalam with śubha madhyāhnaṁ, and say how it differs from the everyday afternoon word.}

Line up the three sister languages' everyday afternoon words and show where all three converge on śubha + madhyāhna.
\end{chapteropening}

\section[śubha madhyāhnaṁ]{\ml{ശുഭ} \ml{മധ്യാഹ്നം} (śubha madhyāhnaṁ) — formal good afternoon}

\label{lesson:ML-C31-shubha-madhyaahnam}



\begin{quote}
Morning broke the pattern entirely. Evening kept the pattern but swapped in an unexpected word. This final greeting brings a different kind of surprise: a word you haven't met, that turns out to connect Malayalam right back to Kannada and Telugu.
\end{quote}

\subsection*{What to know first}
\textbf{\ml{ശുഭ} \ml{മധ്യാഹ്നം}} (\textbf{śubha madhyāhnaṁ}) — \textquotedblleft{}\textbf{good afternoon}\textquotedblright{} — is confirmed, by a directly-fetched source, as a genuinely \textbf{formal} greeting: \textquotedblleft{}\textbf{mainly used in official or ceremonial settings}\textquotedblright{} and \textquotedblleft{}\textbf{not as commonly used in daily conversation}.\textquotedblright{} (A suspiciously precise \textquotedblleft{}fewer than 1\%\textquotedblright{} statistic turned up in an initial search summary for this claim — but it couldn't be traced to either page actually fetched, so it's deliberately left out here as an apparent confabulation; the real, hedged qualitative finding stands on its own.) It's built from \textbf{\ml{ശുഭ}} (ML-C25) plus \textbf{\ml{മധ്യാഹ്നം}} (\emph{madhyāhnaṁ}, \textquotedblleft{}\textbf{noon}\textquotedblright{}) — which Wiktionary confirms is a direct \textbf{synonym} of ML-C17's \textbf{\ml{ഉച്ച}}. Here's the honest surprise: Malayalam had this \textquotedblleft{}middle\textquotedblright{}-based noon-word all along — ML-C17 simply didn't highlight it, teaching \textbf{\ml{ഉച്ച}} (\textquotedblleft{}peak, high\textquotedblright{}) instead. And this greeting is \textbf{not} built on ML-C28's \textbf{\ml{ഉച്ചകഴിഞ്ഞ്}} compound (the actual TIME-AFTERNOON phrase already taught) at all.

\subsection*{Your turn}
Say these aloud:

\begin{itemize}
\raggedright
  \item \textquotedblleft{}śubha madhyāhnaṁ\textquotedblright{} — \textquotedblleft{}good afternoon,\textquotedblright{} a formal, ceremonial-register greeting
  \item the honest surprise — madhyāhnaṁ is a genuine synonym of ucca, Malayalam's own \textquotedblleft{}noon\textquotedblright{} word, just not the one taught first
\end{itemize}

\begin{tcolorbox}[breakable,colback=teal!4,colframe=teal!35!black,title={Before you move on}]
Is \ml{ശുഭ} \ml{മധ്യാഹ്നം} built on ML-C28's \ml{ഉച്ചകഴിഞ്ഞ്}? (\textbf{No} — it's built on \ml{മധ്യാഹ്നം}, a genuine Wiktionary-confirmed synonym of ML-C17's \ml{ഉച്ച}, not the TIME-AFTERNOON compound.) Is the \textquotedblleft{}fewer than 1\% use this\textquotedblright{} statistic something this lesson confirms? (\textbf{No} — it couldn't be traced to any directly-fetched source and is deliberately excluded; only the hedged \textquotedblleft{}formal, not common in daily conversation\textquotedblright{} finding is asserted.)
\end{tcolorbox}
\section[ū]{\ml{ഊ} — one character, met inside words you already say}

\label{lesson:ML-S125-letter-uu}



\begin{quote}
Before the new one: \ml{ഭ} — what does it do?

One character this time. One only. It has not been on a page yet: it arrives with the word on the next one.
\end{quote}

\subsection*{Script you'll notice: \ml{ഊ}}
\textbf{\ml{ഊ}} — \emph{ū}.

It is an \textbf{independent vowel} — the shape this vowel takes when a word \emph{begins} with it, rather than the sign it becomes inside a word.

You already read its short twin. \textbf{\ml{ഉ}} is \emph{u}; \textbf{\ml{ഊ}} is the same vowel held longer, \emph{ū}, and the shape says so — the short one grows a tail.

\begin{itemize}
\raggedright
  \item \textbf{\ml{ഉ}} \emph{u} — the short one, read since the farewell chapter
  \item \textbf{\ml{ഊൺ}} \emph{ūṇ} — a meal: the long one, and the word that will put this letter to work
\end{itemize}

Length is the whole difference, in the sound and on the page.

\subsection*{Writing: \ml{ഊ} — copy what you see}
Put your pen on \ml{ഊ} and follow its line. Copy the shape you can see — slowly, and larger than it is printed.

\begin{quote}
This book does not yet tell you \textbf{where to start the character or which way to travel}. That is a real thing, taught with real variation from school to school, and it is not written down here until it can be written down with a source. Copying what is in front of you needs no such source, and it is how the shape gets into your hand in the meantime.
\end{quote}

\subsection*{Your turn}
\begin{itemize}
\raggedright
  \item \emph{Look:} at these words, and find \ml{ഊ} in the ones that have it
\end{itemize}

\begin{quote}
\ml{ഊൺ}  ·  \ml{ഉ}  ·  \ml{ഊ}
\end{quote}

\begin{itemize}
\raggedright
  \item \emph{Trace it:} \ml{ഊ} three times, saying \emph{ū} as you finish each one
  \item \emph{Look:} back at any page of this chapter and find \ml{ഊ} once more
\end{itemize}

\begin{tcolorbox}[breakable,colback=teal!4,colframe=teal!35!black,title={Before you move on}]
Which character is this — \ml{ഊ}? What sound does it carry? (\emph{\textbf{ū}}.) Which of the two is longer — \ml{ഉ} or \ml{ഊ}? (\ml{ഊ}.)
\end{tcolorbox}
\section[śubha madhyāhnaṁ — three-way convergence]{One greeting, three afternoon strategies}

\label{lesson:ML-C31-afternoon-convergence}



\begin{quote}
Malayalam’s formal greeting uses \textbf{\ml{ശുഭ} + \ml{മധ്യാഹ്നം}}. Kannada and Telugu land on the same phrase—but their ordinary afternoon words took three different routes.
\end{quote}

\subsection*{What to know first}
\noindent
\begin{tabularx}{\linewidth}{@{}>{\raggedright\arraybackslash}X>{\raggedright\arraybackslash}X@{}}
\toprule
\textbf{language} & \textbf{formal afternoon greeting} \\
\midrule
Malayalam & \textbf{\ml{ശുഭ} \ml{മധ്യാഹ്നം}} \\
Kannada & \textbf{\ka{ಶುಭ} \ka{ಮಧ್ಯಾಹ್ನ}} \\
Telugu & \textbf{\te{శుభ} \te{మధ్యాహ్నం}} \\
\bottomrule
\end{tabularx}

All three share the Sanskrit \textbf{śubha + madhyāhna} structure. That shared greeting is striking because the languages’ everyday time words diverge:

\begin{itemize}
\raggedright
  \item \textbf{Telugu} widened its \emph{madhya}-based noon-word so the same word also covers afternoon.
  \item \textbf{Kannada} built a separate Sanskrit tatsama, \textbf{\ka{ಅಪರಾಹ್ನ}} (\emph{aparāhna}, \emph{apara} + \emph{ahna}), rather than widening \emph{madhyāhna}.
  \item \textbf{Malayalam} built \textbf{\ml{ഉച്ചകഴിഞ്ഞ്}}, the unrelated \emph{ucca} “noon” plus “having passed.”
\end{itemize}

So there are three genuinely different everyday strategies, not a simple two-versus-one split. At the formal greeting layer, all three converge on the same borrowed building blocks.

\subsection*{Your turn}
Say these aloud:

\begin{itemize}
\raggedright
  \item the shared greeting structure — \emph{śubha + madhyāhna}
  \item Telugu — widening
  \item Kannada — a separate \emph{aparāhna} word
  \item Malayalam — “noon having passed”
\end{itemize}

\subsection*{Script check — the letters of this chapter}
Cover the romanizations. Find each letter once in the Malayalam above, then say the word it sits in.

\begin{itemize}
\raggedright
  \item \textbf{\ml{ഊൺ}} — \textbf{\ml{ഊ}} \emph{ū}, the long twin of \textbf{\ml{ഉ}}
\end{itemize}

\begin{itemize}
\raggedright
  \item \emph{Look:} at each letter, then find it in the word beside it
\end{itemize}

\begin{tcolorbox}[breakable,colback=teal!4,colframe=teal!35!black,title={Before you move on}]
Do the three languages share one ordinary afternoon-word strategy? (\textbf{No.}) Where do they converge? (\textbf{In the formal greeting, all using śubha + madhyāhna.})
\end{tcolorbox}
